\subsection{Prospective controlled hybrid comparison}
\label{subsec:matched-comparison}
A separate prospective study compared dense--BM25 reciprocal-rank fusion (H) with two-seed citation-graph retrieval (T) on the same 350 queries. Both arms shared the query-only router, evidence formatter, 6,000-token evidence budget, verifier, generator and two evidence-faithfulness judges. This is a controlled implementation comparison, separate from the historical multi-baseline evaluation above. Among the 100 assessable Category~B pairs out of 100 assigned, the T-minus-H composite difference was $0.26885$ (paired percentile-bootstrap 95\% CI $[0.16035, 0.37205]$; two-sided paired randomization $p=0.000010$). T produced 69/100 verifier-triggered refusals on Category~B versus 4/100 for H; this coverage tradeoff prevents interpreting the faithfulness difference as better answer coverage or legal correctness. Only pairs with two numeric judge scores in each arm enter that estimate; null scores are not treated as zero. Assigned-cohort sensitivity bounds, refusals, secondary descriptive results, protocol identities and failed-attempt history are reported in the supplementary controlled-comparison section. Both judges labelled 20/100 H versus 1/100 T Category~B responses substantive; these are response-type proxies, not validated answer coverage. Across all 350 assigned queries per arm, H/T verifier refusals were 172/272 and generator invocations 178/78; both judges called 77/29 responses substantive, while response-type disagreement affected 164/116. On Category~B, generator invocations were 96/31. Invocation and substantive-label counts are not validated correct-answer yield; independent legal-correctness and verifier-error rates remain unestimated. Useful-answer and verifier-error limits are detailed in Supplementary Section~\ref{supp-supp:utility-audit}. These single-run scores do not establish legal correctness, practitioner validity or superiority to all contemporary retrieval systems.
